% Source: fixed 2015 Fujian arts paper, printed page 946.
% From y^2=2px, A=(2,m), |AF|=3 one obtains p=2 and m^2=8.
% The source sketch places A on the upper branch, selecting m=2sqrt(2);
% hence F=(1,0) and the second intersection on the extended AF is B=(1/2,-sqrt(2)).
\begin{tikzpicture}[x=13mm,y=13mm,>=Stealth,
  every node/.style={font=\fontsize{10.5pt}{12.6pt}\selectfont,inner sep=1pt},
  every path/.style={line width=.45pt}]
  \draw[->] (-1.35,0) -- (3.05,0) node[right] {\(x\)};
  \draw[->] (0,-2.55) -- (0,3.25) node[above] {\(y\)};
  % y^2=4x, shown parametrically so the plotted curve matches the derived data.
  \draw[thick,domain=-2.35:3.0,samples=100,smooth,variable=\t]
    plot ({\t*\t/4},{\t});
  \coordinate (O) at (0,0);
  \coordinate (G) at (-1,0);
  \coordinate (F) at (1,0);
  \coordinate (A) at (2,{2*sqrt(2)});
  \coordinate (B) at (.5,{-sqrt(2)});
  \draw (-1.25,{-sqrt(2)/6}) -- (2.25,{13*sqrt(2)/6});
  \draw (-1.25,{sqrt(2)/6}) -- (1.1,{-7*sqrt(2)/5});
  \draw (.15,{-17*sqrt(2)/10}) -- (2.35,{27*sqrt(2)/10});
  \node[below left] at (O) {\(O\)};
  \node[below] at (G) {\(G\)};
  \node[below] at (F) {\(F\)};
  \node[above right,fill=white,inner sep=2pt] at (A) {\(A\)};
  \node[below right] at (B) {\(B\)};
\end{tikzpicture}
